\section{Delphi: qué objeto ético predice realmente la máquina}

La tercera investigación extiende el commonsense a las morally salient situations. Aquí el riesgo aumenta de forma notable: presentar el juicio descriptivo de la mayoría como «moral correcta» reviste el dataset bias de autoridad. Por ello, primero hay que definir el prediction target y solo después hablar de accuracy.

\subsection{La Machine Ethics no es un problema del futuro}

La title slide del paper de Delphi define el sistema como un experiment. La slide siguiente afirma que el despliegue a gran escala de los modelos de lenguaje neuronales exige machine ethics, y cita una advertencia del tipo «la primera ley de la robótica es solo un hombre de paja»: los sistemas reales no se enfrentan a una tabla de reglas morales unificada y enumerable, sino que toman continuamente decisiones implícitas en el filtrado de contenido, la clasificación, las recomendaciones y la interacción.

\lecturefigure{slide-24-delphi-paper.jpg}{Can Machines Learn Morality? The Delphi Experiment}{Diapositiva V024 recuperada del video oficial; Jiang et al., 2021}

\lecturefigure{slide-25-machine-ethics-deployment.jpg}{Los neural systems desplegados ampliamente convierten la machine ethics en una restricción real}{Diapositiva V025 recuperada del video oficial; Stanford CS25 Lecture 16}

\teachervoice{La ponente no estudia machine ethics porque Delphi sea perfecto. Su argumento es justamente el contrario: los sistemas actuales ya toman decisiones con consecuencias éticas; si los investigadores se retiran con el pretexto de que «la moral es demasiado difícil», las reglas predeterminadas y los sesgos existentes seguirán operando.}

\subsection{Las tres Query Interface de Ask Delphi}

La slide de la demo pública muestra la entrada free-form, la respuesta del sistema y un disclaimer visible. La interfaz advierte «research prototype, still work in progress»; esto forma parte de la definición del sistema, no es una nota legal al pie. Le indica al usuario que la salida proviene de una predicción del modelo y que no puede ser el único fundamento de acciones reales, castigos o resoluciones.

\lecturefigure{slide-26-ask-delphi-demo-disclaimer.jpg}{Demo de Ask Delphi y aviso de exención de responsabilidad del prototipo de investigación}{Diapositiva V026 recuperada del video oficial; Stanford CS25 Lecture 16}

La página siguiente coloca lado a lado free-form QA, relative QA y yes/no QA. Free-form evalúa una sola situación, relative compara dos conductas y yes/no responde a un juicio específico. La forma de la tarea cambia el label space y también la superficie de ataque: el relative mode se presta con facilidad a construir comparaciones absurdas entre dos males, y en clase se explicó que esa función se retiró a las pocas horas de hacerse pública.

\lecturefigure{slide-27-delphi-freeform-relative-yesno.jpg}{Interfaces de juicio free-form, relative y yes/no de Delphi}{Diapositiva V027 recuperada del video oficial; Jiang et al., 2021}

\begin{importantbox}{Descriptive Prediction, no Normative Proof}
Lo que Delphi aprende es «cómo suele evaluar la gente, dentro de las poblaciones y situaciones que cubren los datos de entrenamiento», lo que puede escribirse como $p(y\mid s,q)$. Esto no deriva $y$ de axiomas éticos, ni garantiza que la minority view, los principios de justicia o los contextos culturales nuevos queden bien representados.
\end{importantbox}

\subsection{Por qué las Compositional Situations son más difíciles}

La slide de composición parte de «cortar el pasto» y va añadiendo tiempo, lugar, relaciones y condiciones de excepción. Una acción aislada suele tener una evaluación por defecto, pero un context modifier puede invertir el judgment. Al leer la figura, conviene comparar cómo cambia una misma action bajo distintas clause; un sistema que solo memoriza palabras clave pasará por alto condiciones decisivas como «a medianoche», «en horario de trabajo del amigo» o «en una emergencia».

\lecturefigure{slide-28-delphi-compositional-situations.jpg}{Prueba de robustez de Delphi ante situaciones compuestas y modificadores de contexto}{Diapositiva V028 recuperada del video oficial; Jiang et al., 2021}

El par siguiente contrasta conventional scientific knowledge con common life know-how. El modelo sabe que «mezclar cloro con amoniaco es peligroso», pero puede juzgar mal un consejo cotidiano sobre la lactose intolerance. Este contraste muestra que los hechos enciclopédicos, los hechos médicos y las everyday norm provienen de distribuciones de datos distintas; una sola puntuación global oculta las diferencias entre tipos de capacidad.

\lecturefigure{slide-29-conventional-vs-common-life-knowledge.jpg}{Diferencia entre conventional scientific knowledge y common life know-how}{Diapositiva V029 recuperada del video oficial; Stanford CS25 Lecture 16}

\subsection{Evaluation: primero la Task Family, después la puntuación global}

La figura de resultados compara GPT-3 zero-shot, few-shot y 30-shot con Delphi en free-form, yes/no y relative QA. Delphi obtiene valores más altos en las tareas mostradas en clase, pero estos resultados dependen de datos de entrenamiento especializados y de la label definition. Al leer la figura, no debe tratarse la comparación entre un supervised specialist y un generic few-shot model como una architecture-only comparison; la conclusión más razonable es que los targeted moral data cambian de manera notable el judgment behavior.

\lecturefigure{slide-30-delphi-results-vs-gpt3.jpg}{Resultados por tarea de Delphi frente a los prompting baselines de GPT-3}{Diapositiva V030 recuperada del video oficial; Jiang et al., 2021}

Si la situation es $s$, el query format es $q$ y el judgment label es $y$, el objetivo de entrenamiento puede escribirse como
\begin{equation}
\mathcal{L}_{\mathrm{Delphi}}(\theta)
=-\mathbb{E}_{(s,q,y)\sim\mathcal{D}_{\mathrm{norm}}}
\log p_{\theta}(y\mid s,q).
\end{equation}
$\mathcal{D}_{\mathrm{norm}}$ es el conjunto de datos de juicios normativos; esta fórmula muestra directamente que lo que aprende el modelo depende de a quién se muestrea, de cómo se formulan las preguntas y de cómo se combinan las etiquetas.

\subsection{Resumen de la sección}

La task definition de Delphi es el juicio ético descriptivo. Con datos especializados es más estable que el generic prompting, pero la evaluation no puede elevar la accuracy a moral correctness. La siguiente sección debe examinar los datos y el backbone para entender de dónde provienen estos judgment.
